\documentclass[12pt,a4paper]{article}

\usepackage[utf8]{inputenc}
\usepackage[T1]{fontenc}
\usepackage{amsmath,amssymb}
\usepackage{graphicx}
\usepackage{booktabs}
\usepackage{hyperref}
\usepackage[margin=1in]{geometry}
\usepackage{enumitem}
\usepackage{xcolor}
\usepackage{listings}
\usepackage{url}

\hypersetup{
  colorlinks=true,
  linkcolor=blue!70!black,
  citecolor=green!50!black,
  urlcolor=blue!60!black
}

\lstset{
  basicstyle=\ttfamily\small,
  breaklines=true,
  frame=single,
  backgroundcolor=\color{gray!5}
}

\title{\textbf{CryptoGuard: Vision-Language Model-Based Detection of\\
Cryptocurrency Scams in Short-Form Videos}}

\author{
  [Author Name]\\
  [Department], [University]\\
  \texttt{[email]}
}

\date{\today}

\begin{document}

\maketitle

% ===================================================================
\begin{abstract}
Short-form video platforms have become a primary vector for cryptocurrency
scams, combining visual manipulation with high-pressure audio narratives
to deceive viewers. Existing content moderation systems are not specifically
tuned for the multimodal patterns of crypto fraud in video.
We present \textbf{CryptoGuard}, a system that leverages the Image Grid
Vision Language Model (IG-VLM) approach to convert short-form videos into
composite images, combines CLIP-based zero-shot visual scam indicator
scoring with Whisper audio transcription, and feeds this multimodal
evidence into a Qwen3 vision-language model for chain-of-thought scam
detection.
Our pipeline produces structured verdicts with confidence scores and
detailed reasoning, exposed through a web application for interactive
analysis.
We evaluate CryptoGuard on a curated dataset of crypto scam and legitimate
short-form videos, measuring accuracy, precision, recall, and F1 score,
and conduct ablation studies on modality contribution and image grid
configuration.
Results demonstrate that the multimodal approach combining visual, audio,
and CLIP signals achieves superior detection compared to unimodal baselines,
highlighting the potential of VLMs for domain-specific content moderation.
\end{abstract}

\textbf{Keywords:} Vision-Language Models, Cryptocurrency Scam Detection,
Short-Form Video Analysis, Content Moderation, CLIP, Whisper, Multimodal AI

% ===================================================================
\section{Introduction}

Cryptocurrency fraud has surged alongside the growth of digital assets.
The U.S.\ Federal Trade Commission reported that consumers lost over
\$1 billion to cryptocurrency scams between January 2021 and June 2022
alone~\cite{ftc2023}, with social media platforms identified as the
primary channel for these losses. Chainalysis reports similar trends
globally, noting that scam revenue continues to grow year over year~\cite{chainalysis2024}.

Short-form video platforms such as TikTok, YouTube Shorts, and Instagram
Reels have emerged as particularly effective vectors for crypto fraud.
These videos exploit a combination of modalities: visually compelling
fake profit screenshots and luxury lifestyle imagery, urgent spoken
narratives (``only 5 spots left''), and on-screen text with wallet
addresses or QR codes. The multimodal nature of these scams makes them
difficult to detect with text-only or image-only classifiers.

Recent advances in Vision-Language Models (VLMs) have demonstrated
remarkable zero-shot reasoning capabilities across visual and textual
domains~\cite{bordes2024vlm}. Kim et al.~\cite{kim2024imagegrid} showed
that a single VLM can understand video content by converting sampled frames
into an image grid, outperforming dedicated video language models on
9 out of 10 benchmarks. AlDahoul et al.~\cite{aldahoul2024contentmod}
demonstrated that LLMs outperform traditional techniques for content
moderation across text, images, and video.

In this work, we present \textbf{CryptoGuard}, which combines three
key components:
\begin{enumerate}[nosep]
  \item The IG-VLM image grid representation for efficient video-to-image
    conversion~\cite{kim2024imagegrid}.
  \item CLIP~\cite{radford2021clip} zero-shot visual scoring against
    crypto-scam-specific indicator prompts.
  \item Whisper~\cite{radford2023whisper} audio transcription to capture
    spoken scam narratives.
\end{enumerate}
These multimodal signals are fused and analysed by a Qwen3
VLM~\cite{qwen2024} using chain-of-thought prompting to produce
structured scam verdicts.

Our contributions are:
\begin{itemize}[nosep]
  \item A novel application of the IG-VLM image grid technique to
    domain-specific crypto scam detection in short-form video.
  \item A multimodal pipeline combining CLIP visual scoring, Whisper
    transcription, and VLM reasoning.
  \item An evaluation framework with ablation studies on modality
    contribution and grid configuration.
  \item An open-source web application for interactive video analysis.
\end{itemize}

% ===================================================================
\section{Related Work}

\subsection{Video Understanding with Vision-Language Models}

Tang et al.~\cite{tang2023vidllm} provide a comprehensive survey of
approaches for integrating video content with LLMs, categorising methods
into VideoLM pretraining, instruction tuning, LLM-based agents, and
hybrid approaches. A key challenge is encoding temporal information from
potentially hundreds of frames into a format that language models can
process.

Kim et al.~\cite{kim2024imagegrid} address this by proposing Image Grid
VLM (IG-VLM), which converts sampled video frames into a single composite
image arranged in a grid layout. This approach enables direct application
of a frozen, high-performance VLM without any video-specific training.
Their experiments across ten zero-shot video question answering benchmarks
show that IG-VLM surpasses dedicated VideoLMs in 9 out of 10 cases.
The optimal configuration uses 6 frames in a $3 \times 2$ grid with
row-major ordering.

Himakunthala et al.~\cite{himakunthala2023videocot} propose VideoCOT,
which leverages multimodal generative abilities of VLMs to enhance video
reasoning through chain-of-thought processing on keyframes, reducing
computational complexity while improving reasoning quality.

\subsection{Content Moderation with LLMs}

AlDahoul et al.~\cite{aldahoul2024contentmod} evaluate LLM-based content
moderation solutions including OpenAI's moderation model, Llama-Guard-3,
GPT-4o, Gemini 1.5, and Llama-3 across text, images, and videos.
Their results demonstrate that LLMs outperform traditional deep learning
techniques in detecting hate speech, violence, and sexual content, achieving
higher accuracy with lower false positive and false negative rates.
Notably, Gemini 1.5 Pro achieved 95.5\% accuracy for violence detection
in videos.

However, their work focuses on general harmful content categories
(violence, nudity, hate speech) and does not address domain-specific
fraud patterns such as cryptocurrency scams.

\subsection{Misinformation Video Detection}

Bu et al.~\cite{bu2023misinformation} survey misinformation video
detection research, characterising misinformation at three levels:
signals, semantics, and intents. They identify key challenges including
high information heterogeneity across modalities, the blurred distinction
between manipulative and legitimate video editing, and platform-specific
propagation patterns. Their work highlights that domain-specific
misinformation, including financial fraud, remains under-studied.

\subsection{Cryptocurrency Fraud Detection}

Existing work on cryptocurrency fraud focuses primarily on blockchain
transaction analysis using graph neural networks and traditional ML
classifiers operating on on-chain data. Research on detecting crypto
scams through their \emph{promotional content}, particularly in video
format, remains largely unexplored. This gap motivates our work, which
targets the presentation layer of crypto scams rather than their
financial transactions.

% ===================================================================
\section{Methodology}

\subsection{System Architecture}

CryptoGuard processes a short-form video URL through four stages:
(1) video acquisition, (2) multimodal feature extraction, (3) VLM
analysis, and (4) result presentation. Figure~\ref{fig:architecture}
illustrates the pipeline.

\begin{figure}[h]
\centering
\fbox{\parbox{0.9\textwidth}{\centering
\small
\texttt{URL} $\rightarrow$ \texttt{yt-dlp}
$\rightarrow$ \begin{tabular}{l}
  Frame extraction $\rightarrow$ CLIP scoring \\
  Audio extraction $\rightarrow$ Whisper transcript
\end{tabular}
$\rightarrow$ \texttt{Qwen3-VL} $\rightarrow$ \texttt{Verdict}
}}
\caption{CryptoGuard pipeline architecture.}
\label{fig:architecture}
\end{figure}

\subsection{Video Acquisition and Preprocessing}

Videos are downloaded using yt-dlp with a duration constraint
($\leq 180$s) to ensure the input is short-form content. The system
validates the download and extracts metadata for logging.

\subsection{Image Grid Construction}

Following the IG-VLM approach~\cite{kim2024imagegrid}, we uniformly
sample $N = 6$ frames from the video. The video timeline is partitioned
into $N$ equal intervals, and the first frame from each interval is
selected. Frames are resized to $384 \times 384$ pixels and arranged in
a $3 \times 2$ grid (row-major ordering), producing a single
$1152 \times 768$ composite image. This configuration was shown by
Kim et al.\ to be optimal for VLM video understanding, balancing
spatial resolution per frame against temporal coverage.

\subsection{CLIP-Based Visual Scam Indicator Scoring}

We leverage CLIP~\cite{radford2021clip} for zero-shot classification of
visual scam indicators. We define a set of 10 text prompts describing
common crypto scam visual patterns:

\begin{itemize}[nosep]
  \item Fake cryptocurrency profit screenshots
  \item Countdown timers creating urgency
  \item Celebrity endorsement imagery
  \item Luxury lifestyle montages
  \item QR codes or wallet addresses
  \item Fake trading platform interfaces
  \item Text overlays promising giveaways
  \item Before/after wealth testimonials
  \item Fabricated financial charts
  \item Fake bank account balances
\end{itemize}

For each frame, we compute cosine similarity between the CLIP image
embedding and each text prompt embedding. The maximum similarity across
all frames is reported per indicator, producing a score vector
$s \in [0, 1]^{10}$ that quantifies visual scam signal strength.

\subsection{Audio Transcription}

Audio is extracted from the video using ffmpeg and transcribed using
Whisper~\cite{radford2023whisper}. The transcript captures spoken
scam narratives including profit claims, urgency language, and
instructions to send cryptocurrency.

\subsection{VLM Chain-of-Thought Analysis}

The image grid, CLIP indicator scores, and audio transcript are
composed into a structured prompt for the Qwen3 VLM~\cite{qwen2024}.
The system prompt instructs the model to:
\begin{enumerate}[nosep]
  \item Examine the image grid as a temporal sequence of video frames.
  \item Consider the CLIP visual scam indicator scores.
  \item Analyse the audio transcript for scam language patterns.
  \item Reason step-by-step about scam indicators.
  \item Output a structured JSON verdict.
\end{enumerate}

The VLM produces a verdict (\texttt{scam}, \texttt{legitimate}, or
\texttt{uncertain}), a confidence score, chain-of-thought reasoning,
and a list of detected indicators.

% ===================================================================
\section{Implementation}

CryptoGuard is implemented in Python with a FastAPI backend and a
React single-page application frontend. The system supports two
deployment modes:

\begin{description}[nosep]
  \item[Local mode:] Loads Qwen3-VL, CLIP, and Whisper models on a
    local GPU using HuggingFace Transformers.
  \item[API mode:] Calls an OpenAI-compatible API endpoint (e.g., vLLM,
    Ollama) for VLM inference, allowing deployment without local GPU
    resources.
\end{description}

Key implementation details:
\begin{itemize}[nosep]
  \item Video download via \texttt{yt-dlp} with metadata validation.
  \item Frame extraction using OpenCV with uniform temporal sampling.
  \item CLIP encoding via HuggingFace \texttt{CLIPModel} (\texttt{ViT-B/32}).
  \item Audio extraction via ffmpeg, transcription via Whisper \texttt{base} model.
  \item VLM response parsing with JSON extraction and validation.
  \item Lazy model initialization to minimize startup time.
\end{itemize}

The web interface accepts a video URL, displays a loading state during
pipeline execution, and renders the verdict with confidence, reasoning,
CLIP scores, detected indicators, the image grid, and the audio transcript.

% ===================================================================
\section{Evaluation}

\subsection{Dataset Construction}

We construct an evaluation dataset by manually collecting and labelling
short-form video URLs from YouTube Shorts and TikTok. The dataset
contains $N$ samples equally split between crypto scam and legitimate
crypto content. Scam videos are identified by the presence of fake profit
claims, urgency tactics, celebrity impersonation, and suspicious wallet
addresses. Legitimate videos include news coverage, educational content,
and verified project announcements.

\subsection{Metrics}

We evaluate using standard binary classification metrics with ``scam''
as the positive class:
\begin{itemize}[nosep]
  \item \textbf{Accuracy}: overall correct predictions.
  \item \textbf{Precision}: proportion of predicted scams that are true scams.
  \item \textbf{Recall}: proportion of true scams that are detected.
  \item \textbf{F1 Score}: harmonic mean of precision and recall.
  \item \textbf{FPR / FNR}: false positive and false negative rates.
\end{itemize}

\subsection{Ablation Studies}

To measure the contribution of each modality, we conduct ablation
experiments:
\begin{enumerate}[nosep]
  \item \textbf{Visual only}: Image grid + CLIP scores, no transcript.
  \item \textbf{Audio only}: Transcript only, no image grid or CLIP.
  \item \textbf{Full multimodal}: Image grid + CLIP + transcript (our method).
\end{enumerate}

We also vary the image grid configuration:
\begin{itemize}[nosep]
  \item Number of frames: $N \in \{4, 6, 9\}$.
  \item Grid layout: $2 \times 2$, $3 \times 2$, $3 \times 3$.
\end{itemize}

\subsection{Results}

\emph{Results will be populated after evaluation on the curated dataset.
Placeholder tables are provided below.}

\begin{table}[h]
\centering
\caption{Detection performance across modality configurations.}
\label{tab:results}
\begin{tabular}{lccccc}
\toprule
\textbf{Configuration} & \textbf{Acc.} & \textbf{Prec.} & \textbf{Rec.} & \textbf{F1} & \textbf{FNR} \\
\midrule
Visual only        & --  & --  & --  & --  & -- \\
Audio only         & --  & --  & --  & --  & -- \\
Full multimodal    & --  & --  & --  & --  & -- \\
\bottomrule
\end{tabular}
\end{table}

\begin{table}[h]
\centering
\caption{Effect of image grid configuration.}
\label{tab:ablation_grid}
\begin{tabular}{cccc}
\toprule
\textbf{Frames} & \textbf{Grid} & \textbf{Acc.} & \textbf{F1} \\
\midrule
4  & $2 \times 2$ & -- & -- \\
6  & $3 \times 2$ & -- & -- \\
9  & $3 \times 3$ & -- & -- \\
\bottomrule
\end{tabular}
\end{table}

\subsection{Qualitative Analysis}

\emph{Case studies of correct and incorrect classifications will be
presented, analysing which scam indicators the system identifies most
reliably and characterising failure modes.}

% ===================================================================
\section{Discussion}

\paragraph{Strengths.}
The multimodal approach captures scam signals that unimodal methods miss.
CLIP provides interpretable, per-indicator scoring that serves both as
input to the VLM and as explainable evidence to users. The image grid
technique avoids the need for video-specific model training, enabling
rapid deployment.

\paragraph{Limitations.}
The system depends on the VLM's zero-shot reasoning, which can be
sensitive to prompt phrasing and may produce hallucinated indicators.
Short-form video scam tactics evolve rapidly, and the fixed set of CLIP
indicator prompts may not capture novel patterns. Processing time
(download + extraction + inference) is on the order of seconds, which
limits real-time applicability at platform scale.

\paragraph{Ethical Considerations.}
Automated scam detection systems risk false positives that suppress
legitimate content. We recommend human-in-the-loop verification for
enforcement actions. The system processes publicly available videos and
does not collect user data.

% ===================================================================
\section{Conclusion and Future Work}

We presented CryptoGuard, a VLM-based system for detecting cryptocurrency
scams in short-form videos. By adapting the IG-VLM image grid technique
and combining CLIP visual scoring with Whisper transcription, our
multimodal pipeline delivers structured scam verdicts with chain-of-thought
reasoning.

Future work includes: (1) expanding the evaluation dataset for
statistical significance, (2) fine-tuning the VLM on labelled crypto
scam data, (3) incorporating on-screen text extraction (OCR) as an
additional modality, (4) adapting the CLIP indicator prompt set through
prompt optimization, and (5) investigating temporal patterns specific
to scam video editing styles.

% ===================================================================
\bibliographystyle{plain}
\bibliography{references}

\end{document}
